\SourceLeaf{334}%
\[
\begin{aligned}
\log t={}&\log\tau+\frac{Mk^4}{12aa(1-ee)}\mathord{.}rr+\tfrac1{24}M\tau\tau\\
&+\frac M{24k^4}(5ee+(4ee-14e^4)ss+5e^4s^4)\beta\beta\\
\log b={}&\log\beta+\frac{Mk^4}{12aa(1-ee)}\mathord{.}rr+\tfrac18M\tau\tau\\
&-\frac M{24k^4}(2+ee-(8ee-14e^4)ss-9e^4s^4)\beta\beta\\
\log l={}&\log\lambda+\tfrac1{24}M\tau\tau-\frac{(1-ee)M}{24k^4}(1-10eess)\beta\beta
\end{aligned}
\]

Wie man leicht sieht, wird in allen bisher entwickelten Formeln vorausgesetzt, dass die Größen \(t,\tau,b,\beta,l,\lambda\) in Teilen des Radius ausgedrückt sind. Sollen diese Größen in Bogensekunden ausgedrückt und dieselben Bezeichnungen für sie beibehalten werden, so muss man den Formeln für \(\tau,\beta,\lambda\) (Art. 29) noch den Faktor \(\frac1\rho\) hinzufügen. In den Gleichungen 9, 10, 11 und den daraus abgeleiteten logarithmischen Gleichungen muss man dagegen den Gliedern, die \(\tau\tau\) oder \(\beta\beta\) enthalten, noch den Faktor \(\rho\rho\) hinzufügen. Dabei bezeichnet \(\rho\), wie schon in Art. 16 und 19, die Länge eines Bogens von einer Bogensekunde in Teilen des Radius. Behält man nun auch für \(\mu\) die oben verwendete Bedeutung bei, nämlich
\[
\mu=\tfrac1{12}M\rho\rho
\]
und schreibt zur Abkürzung
\[
\begin{aligned}
(1)&=\frac{k}{a\rho}\\
(2)&=\frac{k^3}{a(1-ee)\rho}\\
(3)&=\frac{Mk^4}{12aa(1-ee)}\\
(4)&=\frac{\mu}{2k^4}(5ee+(4ee-14e^4)ss+5e^4s^4)\\
(5)&=\frac{\mu}{2k^4}(2+ee-(8ee-14e^4)ss-9e^4s^4)\\
(6)&=\frac{(1-ee)\mu}{2k^4}(1-10eess)\\
(7)&=\tfrac12\mu
\end{aligned}
\]
so ist unsere Lösung in den folgenden sechs Formeln enthalten
\[
\begin{aligned}
\tau&=(1)r\sin T\operatorname{tang}B\\
\beta&=(2)r\cos T\\
\lambda&=(1)r\sin T\sec B\\
\log t&=\log\tau+(3)rr+(4)\beta\beta+(7)\tau\tau\\
\log b&=\log\beta+(3)rr-(5)\beta\beta+3(7)\tau\tau\\
\log l&=\log\lambda-(6)\beta\beta+(7)\tau\tau
\end{aligned}
\]
